\documentclass[12pt]{article}

\input{~/School/Notes/header.tex}

\title{Домашнее задание по алгебре}
\author{Денисов Илья}
\date{}

\begin{document}
	\maketitle

	% \section*{Задача 1}

	% \subsection*{а)}

	% \noindent
	% \textbf{Условие.} Дан многочлен $f$ над $\mathbb{F}_p$ степени не выше $p-1$ такой, что $f(i) = b_i$, $i = 0, 1, \dots, p-1$. Докажите, что при $0 \leqslant k \leqslant p-2$ коэффициент при $x^{p-1-k}$ равен $-\sum\limits_{i=0}^{p-1} b_i i^k$ (при $k = 0$ считаем $0^0 = 1$).

	% \begin{proof}
	% 	Многочлен степени не выше $p-1$ однозначно определяется значениями в $p$ точках: разность двух таких многочленов имеет $p$ корней и степень не выше $p-1$, значит, равна нулю. Поэтому достаточно предъявить один такой многочлен. Положим
	% 	\[
	% 		F(x) = \sum_{i=0}^{p-1} b_i \bigl( 1 - (x - i)^{p-1} \bigr).
	% 	\]
	% 	Его степень не выше $p-1$. По малой теореме Ферма $(l - i)^{p-1} = 1$ при $l \neq i$, и $(i - i)^{p-1} = 0$. Значит, $F(l) = b_l$, то есть $F = f$.

	% 	Раскроем скобки. Для $0 \leqslant j \leqslant p-1$
	% 	\[
	% 		\binom{p-1}{j} = \prod_{m=1}^{j} \frac{p - m}{m} \equiv \prod_{m=1}^{j} \frac{-m}{m} = (-1)^j \pmod p,
	% 	\]
	% 	поэтому
	% 	\[
	% 		(x - i)^{p-1} = \sum_{j=0}^{p-1} \binom{p-1}{j} x^{p-1-j} (-i)^j = \sum_{j=0}^{p-1} x^{p-1-j} i^j.
	% 	\]
	% 	Слагаемое $1$ в $F$ влияет только на коэффициент при $x^0$, то есть при $k = p-1$, а он нас не интересует. При $0 \leqslant k \leqslant p-2$ коэффициент при $x^{p-1-k}$ в $F$ равен
	% 	\[
	% 		-\sum_{i=0}^{p-1} b_i i^k. \qedhere
	% 	\]
	% \end{proof}

	% \subsection*{б)}

	% \noindent
	% \textbf{Условие.} Многочлены $f, g \in \mathbb{Z}[x]$ задают взаимно обратные перестановки остатков по модулю $17$. Докажите, что если $\deg f = 15$ и старший коэффициент равен $1$, то $\deg g \geqslant 15$. Важно ли условие на старший коэффициент?

	% \begin{proof}
	% 	Здесь $p = 17$, $p - 2 = 15$. Так как $\deg f = 15 \leqslant p-1$, многочлен $f \bmod 17$ сам является интерполяционным многочленом функции $i \mapsto f(i)$. Применим а) при $k = 1$ ($b_i = f(i)$): коэффициент при $x^{15}$ равен $-\sum\limits_i i f(i)$. По условию он равен $1$, то есть
	% 	\[
	% 		\sum_{i=0}^{16} i f(i) \equiv -1 \pmod{17}.
	% 	\]

	% 	Пусть $\hat g$~--- интерполяционный многочлен функции $j \mapsto g(j) \bmod 17$ степени не выше $16$. Он получается из $g$ приведением по модулю $x^{17} - x$ и модулю $17$, а при таком приведении степень не растёт, поэтому $\deg \hat g \leqslant \deg g$. По а) при $k = 1$ коэффициент $\hat g$ при $x^{15}$ равен $-\sum\limits_j j\, g(j)$. Так как $f$ и $g$ взаимно обратны, $j = f(i)$ пробегает все остатки и $g(f(i)) = i$:
	% 	\[
	% 		-\sum_{j} j\, g(j) = -\sum_{i} f(i)\, i = 1 \neq 0.
	% 	\]
	% 	Значит, $\deg \hat g \geqslant 15$, откуда $\deg g \geqslant \deg \hat g \geqslant 15$.
	% \end{proof}

	% \noindent
	% \textbf{Про старший коэффициент.} Само значение $1$ не важно: если старший коэффициент равен $a \not\equiv 0 \pmod{17}$, то коэффициент $\hat g$ при $x^{15}$ тоже равен $a$, и рассуждение то же. Важно лишь, чтобы $a$ не делилось на $17$. Иначе утверждение неверно: например, $f = x + 17x^{15}$ и $g = x$ задают одну и ту же (тождественную) перестановку, $\deg f = 15$, но $\deg g = 1$.

	% %==================================================
	% \section*{Задача 2}

	% \noindent
	% \textbf{Условие.} Линейная комбинация $S$ чисел $\cos 20^\circ$, $\cos 40^\circ$, $\cos 80^\circ$ с целыми коэффициентами оказалась корнем из целого числа. Докажите, что $S$~--- целое число.

	% \begin{proof}
	% 	Пусть $a = 2\cos 40^\circ$. Из $2\cos 3\varphi = (2\cos\varphi)^3 - 3 \cdot 2\cos\varphi$ при $\varphi = 40^\circ$ получаем $a^3 - 3a = 2\cos 120^\circ = -1$, то есть $a$~--- корень $x^3 - 3x + 1$. У этого многочлена нет рациональных корней ($\pm 1$ не подходят), значит, он неприводим над $\mathbb{Q}$ и $[\mathbb{Q}(a) : \mathbb{Q}] = 3$.

	% 	Выразим остальные числа через $a$. Так как $2\cos 80^\circ = (2\cos 40^\circ)^2 - 2$, имеем $b := 2\cos 80^\circ = a^2 - 2$. Далее $2\cos 160^\circ = -2\cos 20^\circ$ и $a + b + 2\cos 160^\circ = 0$, поскольку $\cos 40^\circ + \cos 80^\circ + \cos 160^\circ = 0$ (сумма корней $x^3 - 3x + 1$ равна нулю). Значит,
	% 	\[
	% 		2\cos 20^\circ = a + b = a^2 + a - 2.
	% 	\]
	% 	Пусть $S = m \cos 20^\circ + n \cos 40^\circ + l \cos 80^\circ$, $m, n, l \in \mathbb{Z}$. Тогда
	% 	\[
	% 		2S = m (a^2 + a - 2) + n a + l (a^2 - 2) = (m + l) a^2 + (m + n) a - 2(m + l),
	% 	\]
	% 	и $2S \in \mathbb{Q}(a)$.

	% 	По условию $S^2 = N \in \mathbb{Z}$. Предположим, что $2S \notin \mathbb{Q}$. Тогда $2S$~--- корень $x^2 - 4N$, неприводимого над $\mathbb{Q}$, и $[\mathbb{Q}(2S) : \mathbb{Q}] = 2$. Но $\mathbb{Q}(2S) \subset \mathbb{Q}(a)$, и число $2$ должно делить $[\mathbb{Q}(a) : \mathbb{Q}] = 3$. Противоречие. Итак, $2S \in \mathbb{Q}$.

	% 	Числа $1, a, a^2$ линейно независимы над $\mathbb{Q}$, поэтому равенство $2S = (m+l)a^2 + (m+n)a - 2(m+l) \in \mathbb{Q}$ возможно только при $m + l = 0$ и $m + n = 0$. Тогда $2S = 0$, то есть $S = 0 \in \mathbb{Z}$.
	% \end{proof}

	%==================================================
	\section*{Задача 3}

	\noindent
	\textbf{Условие.} Сколько решений в рациональных числах может иметь уравнение $x^3 + kx^2 - (k+3)x + 1 = 0$, где $k$~--- вещественный параметр?

	\noindent
	\textbf{Ответ.} Только $0$ или $3$ (оба случая реализуются).

	\begin{proof}
		Обозначим левую часть $f(x)$. Подстановка $x = 0$ даёт $f(0) = 1 \neq 0$, а $x = 1$ даёт $f(1) = -1 \neq 0$. Поэтому для рационального корня $x \ne 0, 1$ из уравнения
		\[
			k(x^2 - x) = 3x - 1 - x^3 \quad \Longrightarrow \quad k = \frac{3x - 1 - x^3}{x^2 - x} \in \mathbb{Q}.
		\]
		Значит, при иррациональном $k$ рациональных корней нет. Пусть теперь $k \in \mathbb{Q}$ и $x_1 \in \mathbb{Q}$~--- корень.

		Дискриминант $D = (x_1 - x_2)^2(x_1 - x_3)^2(x_2 - x_3)^2$:
		\[
			D = p^2q^2 - 4p^3r - 4q^3 + 18pqr - 27r^2 = k^4 + 6k^3 + 27k^2 + 54k + 81 = (k^2 + 3k + 9)^2.
		\]
		Так как $k^2 + 3k + 9 > 0$ (дискриминант этого трёхчлена равен $-27$), $D > 0$ и $D$~--- квадрат рационального числа. В частности, корни попарно различны.

		Разделим $f$ на $x - x_1$: $f(x) = (x - x_1)(x^2 + \alpha x + \beta)$, $\alpha, \beta \in \mathbb{Q}$. Пусть $d = \alpha^2 - 4\beta = (x_2 - x_3)^2$. Тогда $(x_1 - x_2)(x_1 - x_3) = f'(x_1) \in \mathbb{Q}$, $f'(x_1) \neq 0$ (корни простые), и
		\[
			D = f'(x_1)^2 \cdot d \quad \Longrightarrow \quad d = \frac{D}{f'(x_1)^2}.
		\]
		Здесь $D$~--- квадрат рационального числа, значит, и $d$~--- квадрат рационального числа. Тогда $x_{2,3} = \dfrac{-\alpha \pm \sqrt d}{2} \in \mathbb{Q}$.

		Итак, если есть хотя бы один рациональный корень, то их три; иначе ни одного. Случаи $1$ и $2$ невозможны.

		\emph{Оба случая реализуются.} При $k = 0$ получаем $x^3 - 3x + 1$~--- он неприводим, корней нет. При $k = -\tfrac32$ имеем $f(x) = (x - 2)\left(x^2 + \tfrac12 x - \tfrac12\right)$, корни $2$, $\tfrac12$, $-1$.
	\end{proof}

\end{document}